%% =====================================================================
%%  T-16-04  The Unargue
%%  -------------------------------------------------------------------
%%  One idea per file. Core is mandatory. Slots in canonical order.
%%  Max 700 words. No layout commands. No filler. New source material
%%  for this entry goes in sources/B1/T-16-04.tex -- never by
%%  rewriting this file (docs/RULES.md R0.1).
%%  =====================================================================
%% @kind: topic
%% @id: T-16-04
%% @title: The Unargue
%% @subtitle: Winning by making agreement the target's own comfortable exit
%% @chapter: c16
%% @order: 40
%% @status: stable
%% @sources: B1
%% @updated: 2026-09-19

\topic{T-16-04}{The Unargue}{Winning by making agreement the target's own comfortable exit}

\begin{core}
An argument yields an angry adversary, more convinced than ever. B1's counter-technique --- the Unargue --- is argument run backwards: never dispute, never defend; build agreement around the disagreement until the target's concession arrives as their own comfortable exit. Six steps, street-tested against one benchmark only: it has to work.
\end{core}
\begin{mechanism}
The technique's engine is ego-management. A person's refusal is rarely about the proposition; it is armour for the self, so every step is designed to keep the armour off. (1) The soulmate tactic: find and dwell on real areas of agreement --- no pretence, since anyone over an IQ of fifty spots a phony --- until you are a friend with one issue left. (2) Listen, listen, listen: hear out the objection fully; separate the token misgivings your momentum can rout from the one or two major objections that actually block you. (3) Agree with the feelings, not the case: meet hostility with \emph{I know how you feel} and a self-deprecating laugh --- the ego is massaged aside without conceding an inch of substance. (4) Point out areas of agreement and appeal to common values: we both want this campaign to make you money. (5) State your case through the nine-out-of-ten device: \emph{you'd usually be right, but this situation is unusual} --- the objection is honoured and set aside in one motion, and as with a lost battle, the rout of minor objections follows. (6) Save face, throughout: concede a minor point to start a pattern of concession, offer a compromise so the target can yield with their head up, and never let the word \emph{wrong} attach to them.
\end{mechanism}
\begin{conditions}
Strongest on the persuadable counterpart in a one-to-one setting with rapport time --- sales, supervision, family negotiation. It decays under audience pressure (face is scarcer in public), against the genuinely immovable (T-16-02's triage governs who gets the effort), and where the disagreement is structural rather than sentimental --- no amount of ego-salve moves a budget line.
\end{conditions}
\begin{application}
  \item Before the meeting, bank your agreements: list what you genuinely share, and open there.
  \item Run step 2 longer than comfort. The target talking is the technique working.
  \item Classify objections live: token or major. Fight only the majors; let momentum mop up the rest.
  \item Retaliate at nothing. Insults get \emph{I know how you feel} and a laugh at your own expense.
  \item Buy the concession with face: your minor apology, your visible compromise, their head held high.
\end{application}
\begin{feedback}
  \item Landing: the target starts arguing your side of the case to themselves. The exit has been found.
  \item Landing: they accept while insisting they have not really changed their mind. That is the face-saving working, not failing.
  \item Failing: you defended yourself once. The armour is back on and the clock has reset.
  \item Failing: they agree to everything and do nothing. The technique moved the mood, not the incentive; check whether the blocker was ever sentimental.
\end{feedback}
\begin{failure}
The method's honest limit: it produces consent, not conversion --- remove the face-saving scaffolding and the target may quietly re-erect the objection. Its darker reading is the source's own: the technique is calibrated to leave the target satisfied while carrying your purpose, which is exactly what \emph{against his will and make him like it} admits. Practitioners should know which side of that sentence they are standing on.
\end{failure}
\begin{countermeasures}
  \item Notice the pattern run on you: warmth, deep listening, your feelings validated, a concession extracted with your dignity intact. Ask what you just agreed to before you leave the room.
  \item Distinguish token objections from real ones on your side too --- the operator's triage works because targets offer their real objection last, wrapped in a decoy (T-05-01).
  \item Buy time before yielding: \emph{let me think about it overnight} defeats the nine-out-of-ten device, which needs its momentum.
  \item If someone always makes you feel wise and comfortable right before you concede, price the comfort. It is a tool, and it has a user (T-13-02).
\end{countermeasures}

\sources{B1}
\seesources
\seealso{T-16-02, T-16-03, T-05-01}
